\chapter{Dockerfile: сборка образов}

Dockerfile~--- текстовое описание того, как собрать образ. Это самый
важный файл в проекте: от него зависят скорость сборки, размер образа,
воспроизводимость и безопасность. В главе разобраны все основные
инструкции, разница между \cmd{CMD} и \cmd{ENTRYPOINT}, порядок
инструкций и кэш слоёв, многоэтапная сборка и современные возможности
BuildKit.

\section{Как работает сборка}

Команда \cmd{docker build} отправляет демону \emph{контекст сборки}~---
каталог, в котором лежит Dockerfile,~--- и последовательно выполняет
инструкции. Каждая инструкция, меняющая файловую систему, создаёт новый
слой поверх предыдущего. Результат~--- образ.

\begin{code}[language=bash]
$ mkdir hello-docker && cd hello-docker
$ cat > Dockerfile <<'EOF'
FROM alpine:3.21
RUN echo "Привет из контейнера" > /hello.txt
CMD ["cat", "/hello.txt"]
EOF
$ docker build -t hello:1.0 .
$ docker run --rm hello:1.0
Привет из контейнера
\end{code}

\begin{warnbox}
Точка в конце команды~--- это контекст сборки, а не «просто так».
\cmd{docker build -t app .} отправляет демону всё содержимое текущего
каталога, кроме перечисленного в \file{.dockerignore}. Сборка
из корня домашнего каталога может занять минуты и отправить гигабайты
мусора.
\end{warnbox}

\begin{defbox}
\textbf{Контекст сборки}~--- набор файлов, доступных инструкциям
\cmd{COPY} и \cmd{ADD}. Инструкции \cmd{RUN} выполняются внутри
контейнера-сборщика и не видят файлы хоста напрямую.
\end{defbox}

\section{Первый настоящий Dockerfile}

Соберём небольшой веб-сервис на Python и Flask.

\begin{code}[language=bash]
$ tree
.
├── Dockerfile
├── app.py
└── requirements.txt
\end{code}

\begin{code}[language=python]
from flask import Flask

app = Flask(__name__)

@app.get("/")
def index():
    return {"status": "ok", "service": "demo"}

@app.get("/health")
def health():
    return {"healthy": True}

if __name__ == "__main__":
    app.run(host="0.0.0.0", port=8000)
\end{code}

\begin{code}
flask==3.1.0
\end{code}

\begin{code}[language=dockerfile]
FROM python:3.13-slim

ENV PYTHONDONTWRITEBYTECODE=1 \
    PYTHONUNBUFFERED=1

WORKDIR /app

COPY requirements.txt .
RUN pip install --no-cache-dir -r requirements.txt

COPY app.py .

EXPOSE 8000
USER 10001

HEALTHCHECK --interval=30s --timeout=3s --start-period=5s --retries=3 \
  CMD python -c "import urllib.request; urllib.request.urlopen('http://127.0.0.1:8000/health')" || exit 1

CMD ["python", "app.py"]
\end{code}

Сборка и запуск:

\begin{code}[language=bash]
$ docker build -t demo-api:1.0 .
$ docker run -d --name api -p 8000:8000 demo-api:1.0
$ curl http://localhost:8000/
{"service":"demo","status":"ok"}
$ docker ps --filter name=api --format '{{.Status}}'
Up 12 seconds (healthy)
\end{code}

Обратите внимание на порядок: сначала копируется только
\file{requirements.txt}, затем устанавливаются зависимости и только
потом~--- код. При изменении \file{app.py} Docker переиспользует слой
с зависимостями и не выполняет \cmd{pip install} заново.

\section{Инструкции Dockerfile}

\subsection{FROM}

Задаёт базовый образ. Должна быть первой инструкцией (кроме \cmd{ARG}
до неё).

\begin{code}[language=dockerfile]
FROM python:3.13-slim
FROM node:22-alpine AS build
FROM scratch
\end{code}

\begin{itemize}
  \item \cmd{scratch}~--- пустой образ. Используется для статических
        бинарников и минимальных образов.
  \item \cmd{AS имя}~--- имя этапа для многоэтапной сборки.
  \item Всегда фиксируйте версию: \cmd{python:3.13-slim} лучше, чем
        \cmd{python:latest}.
\end{itemize}

\subsection{RUN}

Выполняет команду в новом слое во время сборки.

\begin{code}[language=dockerfile]
RUN apt-get update && apt-get install -y --no-install-recommends \
      curl ca-certificates \
    && rm -rf /var/lib/apt/lists/*
RUN pip install --no-cache-dir -r requirements.txt
\end{code}

\begin{warnbox}
Каждый \cmd{RUN} создаёт слой. Раздельные команды
\cmd{RUN apt-get update} и \cmd{RUN apt-get install}~--- плохая идея:
кэш обновится, а слой установки останется старым. Объединяйте связанные
команды через \cmd{\&\&} и удаляйте временные файлы в том же слое.
\end{warnbox}

\subsection{COPY и ADD}

Копируют файлы из контекста в образ.

\begin{code}[language=dockerfile]
COPY package.json package-lock.json ./
COPY --chown=app:app ./src /app/src
COPY --chmod=755 entrypoint.sh /usr/local/bin/entrypoint.sh
\end{code}

\begin{itemize}
  \item \cmd{COPY}~--- предсказуемое копирование. Используйте его
        по умолчанию.
  \item \cmd{ADD} умеет распаковывать локальные tar-архивы и скачивать
        URL. Скачивание по URL лучше заменить на \cmd{RUN curl},
        чтобы контролировать проверку контрольной суммы; распаковку
        архива~--- на \cmd{COPY} + \cmd{RUN tar}.
\end{itemize}

\subsection{WORKDIR}

Задаёт рабочий каталог для последующих инструкций и создаёт его,
если он не существует.

\begin{code}[language=dockerfile]
WORKDIR /app
\end{code}

Не используйте \cmd{RUN cd /app}: эта команда действует только внутри
своего слоя и не влияет на следующие инструкции.

\subsection{ENV и ARG}

Обе инструкции задают переменные, но с разными целями.

\begin{center}
\small
\begin{tabularx}{\textwidth}{@{}lXX@{}}
\toprule
 & \textbf{ARG} & \textbf{ENV} \\
\midrule
Когда доступна & только во время сборки & при сборке и в работающем
контейнере \\
Сохраняется в образе & нет (кроме случаев, когда попала в слой) &
да, в конфигурации образа \\
Переопределение & \flag{--build-arg} & \flag{-e} или \cmd{environment}
в Compose \\
Типичное применение & версии, флаги сборки & конфигурация приложения \\
\bottomrule
\end{tabularx}
\end{center}

\begin{code}[language=dockerfile]
ARG PYTHON_VERSION=3.13
FROM python:${PYTHON_VERSION}-slim

ARG APP_VERSION=dev
ENV APP_VERSION=${APP_VERSION} \
    TZ=Europe/Moscow
\end{code}

\begin{warnbox}
\cmd{ARG} не подходит для секретов. Значение аргумента попадает
в историю сборки и может быть извлечено из образа. Для паролей
и токенов используйте \flag{--secret} (раздел~6.8) или переменные
окружения при запуске.
\end{warnbox}

\subsection{EXPOSE}

Документирует порт, который слушает приложение. \cmd{EXPOSE} не
публикует порт: для этого нужен \flag{-p} при запуске. Инструкция
полезна как подсказка людям и инструментам.

\begin{code}[language=dockerfile]
EXPOSE 8000
EXPOSE 8000/tcp 9090/udp
\end{code}

\subsection{USER}

Задаёт пользователя, от имени которого работает контейнер.

\begin{code}[language=dockerfile]
RUN useradd --create-home --uid 10001 app
USER app
\end{code}

В Debian/Ubuntu-образах удобно \cmd{useradd}; в Alpine~---
\cmd{adduser -D -u 10001 app}. Можно указать числовой UID без создания
пользователя: \cmd{USER 10001}. Kubernetes и некоторые системы
безопасности предпочитают именно числовые UID.

\subsection{CMD и ENTRYPOINT}

Это самая частая тема ошибок. \cmd{ENTRYPOINT} задаёт исполняемый файл,
\cmd{CMD}~--- аргументы по умолчанию (или команду, если
\cmd{ENTRYPOINT} нет).

\begin{center}
\small
\begin{tabularx}{\textwidth}{@{}lXX@{}}
\toprule
\textbf{В Dockerfile} & \textbf{Что выполнится} & \textbf{\cmd{docker run img ARG}} \\
\midrule
только \cmd{CMD} & \cmd{CMD} целиком & заменяет \cmd{CMD} \\
только \cmd{ENTRYPOINT} & \cmd{ENTRYPOINT} & добавляется аргументом \\
оба & \cmd{ENTRYPOINT} + \cmd{CMD} & заменяет \cmd{CMD}, \cmd{ENTRYPOINT}
остаётся \\
ни одного & ошибка при запуске &~--- \\
\bottomrule
\end{tabularx}
\end{center}

\begin{code}[language=dockerfile]
# Вариант 1: ENTRYPOINT + аргументы по умолчанию
ENTRYPOINT ["python", "app.py"]
CMD ["--port", "8000"]

# Вариант 2: только CMD, легко переопределить целиком
CMD ["python", "app.py"]
\end{code}

\begin{code}[language=bash]
$ docker run --rm demo-api:1.0 --port 9000      # аргументы к ENTRYPOINT
$ docker run --rm --entrypoint sh demo-api:1.0  # временно заменить ENTRYPOINT
\end{code}

\begin{warnbox}
Используйте \emph{exec-форму} (\cmd{["python", "app.py"]}), а не
shell-форму (\cmd{python app.py}). В shell-форме команда запускается
через \cmd{/bin/sh -c}, и сигнал \cmd{SIGTERM} не доходит до приложения:
\cmd{docker stop} будет ждать таймаут и затем убивать процесс.
Shell-форма нужна только там, где действительно требуется подстановка
переменных или конвейер.
\end{warnbox}

Если приложению нужны и подстановки, и корректные сигналы, используйте
entrypoint-скрипт с \cmd{exec}:

\begin{code}[language=bash]
#!/bin/sh
set -e
echo "Запуск с PORT=$PORT"
exec "$@"
\end{code}

\begin{code}[language=dockerfile]
COPY entrypoint.sh /usr/local/bin/entrypoint.sh
RUN chmod +x /usr/local/bin/entrypoint.sh
ENTRYPOINT ["entrypoint.sh"]
CMD ["python", "app.py"]
\end{code}

\subsection{HEALTHCHECK}

Позволяет Docker и Compose отличать «процесс жив» от «сервис отвечает».

\begin{code}[language=dockerfile]
HEALTHCHECK --interval=30s --timeout=3s --start-period=10s --retries=3 \
  CMD curl -fsS http://127.0.0.1:8000/health || exit 1
\end{code}

Параметры: \cmd{--interval}~--- как часто проверять, \cmd{--timeout}~
--- сколько ждать ответа, \cmd{--start-period}~--- сколько дать
приложению на старт, \cmd{--retries}~--- сколько неудач подряд считать
болезнью. Проверка должна быть быстрой и дешёвой; не стоит ходить
во внешние сервисы.

\begin{notebox}
HEALTHCHECK не перезапускает контейнер сам по себе. Он влияет на статус
(\cmd{healthy}/\cmd{unhealthy}), а решения принимают Compose
(\cmd{depends\_on: condition: service\_healthy}), Swarm или Kubernetes.
\end{notebox}

\subsection{Остальные инструкции}

\begin{description}[style=nextline]
  \item[LABEL] Метаданные образа: автор, версия, ссылка на репозиторий.
        \cmd{org.opencontainers.\-image.\-version="1.0"}.
  \item[VOLUME] Объявляет точку монтирования. Данные, записанные в неё
        при сборке, теряются: том создаётся при запуске. Используйте
        осознанно, предпочитая тома, создаваемые при запуске.
  \item[STOPSIGNAL] Сигнал, который Docker посылает при остановке.
        По умолчанию \cmd{SIGTERM}.
  \item[SHELL] Меняет оболочку для shell-формы инструкций.
        \cmd{SHELL ["/bin/bash", "-o", "pipefail", "-c"]}.
  \item[ONBUILD] Откладывает инструкцию до сборки дочернего образа.
        Встречается в старых базовых образах; в новом коде избегайте.
\end{description}

\section{Кэш слоёв}

Docker кэширует каждый слой. Инструкция выполняется заново, если
изменилась она сама, её аргументы или любой из предыдущих слоёв.
Поэтому порядок инструкций определяет скорость сборки.

\begin{code}[language=dockerfile]
# Плохо: любое изменение кода сбрасывает установку зависимостей
COPY . .
RUN npm ci

# Хорошо: сначала манифесты, потом код
COPY package.json package-lock.json ./
RUN npm ci
COPY . .
\end{code}

\begin{code}[language=bash]
$ docker build -t app .
[+] Building 1.2s (12/12) FINISHED
 => [2/5] COPY package.json package-lock.json ./
 => [3/5] RUN npm ci
 => CACHED [4/5] COPY . .
\end{code}

Полезные приёмы:

\begin{itemize}
  \item \flag{--no-cache}~--- собрать всё заново (для проверки
        воспроизводимости).
  \item \flag{--progress=plain}~--- подробный вывод, удобно при отладке.
  \item \flag{--pull}~--- обновить базовый образ перед сборкой.
  \item \cmd{docker builder prune}~--- очистить кэш сборки.
\end{itemize}

\begin{tipbox}
В CI кэш можно переносить между сборками через реестр:
\cmd{docker buildx build --cache-from type=registry,ref=...}. Тогда
зависимости не устанавливаются заново на каждом запуске.
\end{tipbox}

\section{.dockerignore}

Файл \file{.dockerignore} исключает лишнее из контекста сборки. Он
ускоряет сборку, уменьшает риск утечки секретов и предотвращает
случайное копирование \cmd{node\_modules} или \file{.git}.

\begin{code}
.git
.gitignore
node_modules
__pycache__
*.pyc
.env
.env.*
*.log
dist
coverage
Dockerfile
docker-compose*.yml
README.md
\end{code}

\begin{warnbox}
Без \file{.dockerignore} команда \cmd{COPY . .} скопирует в образ
\file{.env} с паролями, историю Git и локальные артефакты. Секреты
в образе остаются в слоях даже после удаления файла в следующей
инструкции.
\end{warnbox}

\section{Многоэтапная сборка}

Многоэтапная сборка (\emph{multi-stage build}) позволяет использовать
тяжёлые инструменты на этапе компиляции и оставлять в финальном образе
только результат.

\begin{figure}[htbp]
\centering
\begin{tikzpicture}[
  font=\sffamily\scriptsize,
  stage/.style={draw=accent!55!black, fill=accent!7, rounded corners=2pt,
    minimum width=4.6cm, minimum height=1.5cm, align=left, inner sep=6pt},
  final/.style={draw=good!60!black, fill=good!8, rounded corners=2pt,
    minimum width=4.6cm, minimum height=1.5cm, align=left, inner sep=6pt},
  arr/.style={-Latex, thick, accent!70},
]
  \node[stage] (build) at (0,0)
    {\textbf{этап build}\\[2pt]\cmd{golang:1.25-alpine}\\компилятор, исходники\\\cmd{go build -o /out/app}};
  \node[final, right=2.2cm of build] (final)
    {\textbf{финальный этап}\\[2pt]\cmd{scratch}\\только бинарник\\\cmd{COPY --from=build}};
  \draw[arr] (build) -- node[above, text=muted] {артефакт} (final);
  \node[below=2pt of build, text=muted] {$\approx$ 850 МБ};
  \node[below=2pt of final, text=muted] {$\approx$ 8 МБ};
\end{tikzpicture}
\caption{Многоэтапная сборка: инструменты остаются в первом этапе,
в финальный образ попадает только результат}\label{fig:multistage}
\end{figure}

\begin{code}[language=dockerfile]
FROM golang:1.25-alpine AS build
WORKDIR /src
COPY go.mod go.sum ./
RUN go mod download
COPY . .
RUN CGO_ENABLED=0 go build -trimpath -ldflags="-s -w" -o /out/app ./cmd/app

FROM scratch
COPY --from=build /out/app /app
USER 10001
ENTRYPOINT ["/app"]
\end{code}

\begin{code}[language=bash]
$ docker build -t app:1.0 .
$ docker image ls app
REPOSITORY   TAG   IMAGE ID       CREATED         SIZE
app          1.0   c3d4e5f6a7b8   3 seconds ago   8.1MB
\end{code}

Для Node.js и фронтенда финальным этапом часто служит nginx:

\begin{code}[language=dockerfile]
FROM node:22-alpine AS build
WORKDIR /app
COPY package*.json ./
RUN npm ci
COPY . .
RUN npm run build

FROM nginx:1.27-alpine
COPY --from=build /app/dist /usr/share/nginx/html
\end{code}

\begin{tipbox}
Можно остановиться на промежуточном этапе и зайти в него:
\cmd{docker build --target build -t app:build .}, затем
\cmd{docker run --rm -it app:build sh}. Это удобно, когда сборка падает
и нужно посмотреть, что получилось.
\end{tipbox}

\section{Возможности BuildKit}

Современный Docker по умолчанию использует BuildKit~--- сборщик
с расширенными возможностями. Некоторые из них особенно полезны.

\subsection{Кэш пакетных менеджеров}

\begin{code}[language=dockerfile]
RUN --mount=type=cache,target=/root/.cache/pip \
    pip install -r requirements.txt

RUN --mount=type=cache,target=/root/.npm \
    npm ci
\end{code}

Кэш сохраняется между сборками и не попадает в образ.

\subsection{Секреты при сборке}

\begin{code}[language=dockerfile]
RUN --mount=type=secret,id=npmrc,target=/root/.npmrc \
    npm ci
\end{code}

\begin{code}[language=bash]
$ docker build --secret id=npmrc,src=$HOME/.npmrc -t app .
\end{code}

Секрет доступен только на время выполнения инструкции и не остаётся
в слоях.

\subsection{Here-doc и COPY --link}

\begin{code}[language=dockerfile]
RUN <<EOF
apt-get update
apt-get install -y --no-install-recommends curl
rm -rf /var/lib/apt/lists/*
EOF

COPY --link --chown=10001:10001 . /app
\end{code}

Here-doc избавляет от длинных цепочек \cmd{\&\&}, а \cmd{COPY --link}
улучшает переиспользование слоёв при пересборке.

\subsection{Мультиархитектурная сборка}

\begin{code}[language=bash]
$ docker buildx build \
    --platform linux/amd64,linux/arm64 \
    --tag registry.example.com/team/app:1.0 \
    --push .
\end{code}

Результат~--- один тег, указывающий на образы для двух архитектур.
Требуется драйвер \cmd{docker-container} или удалённый builder;
локальные образы при \flag{--push} не загружаются автоматически.

\section{Типичные ошибки}

\begin{center}
\small
\begin{tabularx}{\textwidth}{@{}XX@{}}
\toprule
\textbf{Ошибка} & \textbf{Последствие} \\
\midrule
\cmd{COPY . .} до установки зависимостей & кэш сбрасывается при каждом
изменении кода \\
\cmd{apt-get update} без \cmd{install} в том же \cmd{RUN} & устаревший
индекс пакетов, ошибки 404 \\
секреты в \cmd{ARG} или \cmd{ENV} & утечка в историю образа \\
\cmd{latest} в базовом образе & невоспроизводимые сборки \\
запуск от \cmd{root} без необходимости & лишние права при компрометации \\
установка компиляторов в финальный образ & образ в разы больше нужного \\
\cmd{ADD} для локальных файлов & неявная распаковка и неожиданное
поведение \\
нет \file{.dockerignore} & гигабайты контекста и утечка \file{.env} \\
shell-форма \cmd{CMD} & сигналы не доходят до приложения \\
\bottomrule
\end{tabularx}
\end{center}

\begin{summarybox}
\begin{itemize}
  \item Dockerfile выполняется последовательно; каждая инструкция,
        меняющая файлы, создаёт слой.
  \item Сначала копируйте манифесты зависимостей, потом код~---
        так работает кэш.
  \item \cmd{CMD}~--- аргументы по умолчанию, \cmd{ENTRYPOINT}~---
        исполняемый файл; используйте exec-форму.
  \item \file{.dockerignore} обязателен; секреты не должны попадать
        в контекст и в \cmd{ARG}.
  \item Многоэтапная сборка оставляет в образе только результат,
        а BuildKit даёт кэш, секреты и мультиархитектурные сборки.
\end{itemize}
\end{summarybox}

\section*{Вопросы для самопроверки}
\begin{selfcheck}
  \item Чем \cmd{COPY} отличается от \cmd{ADD} и что предпочтительнее?
  \item Чем \cmd{ARG} отличается от \cmd{ENV}?
  \item Почему \cmd{CMD ["python", "app.py"]} лучше \cmd{CMD python app.py}?
  \item Почему зависимости устанавливают до копирования всего кода?
  \item Что произойдёт с файлами, записанными в каталог, объявленный
        \cmd{VOLUME}, во время сборки?
  \item Зачем нужна многоэтапная сборка?
  \item Как передать секрет в \cmd{RUN} и не оставить его в образе?
\end{selfcheck}
